\chapter{The Nature of Virtual Universes}

\section{Common Sense}
  
Particles, the flow of time—everything, in fact—emerge as patterns within white noise.
That white noise is abstract information.
Without any physical interpretation attached to it—without apples, bananas, or electrons—such structures are purely abstract.
They are not made of matter. They cannot be touched, weighed, or photographed. But they surely are! Mathematics does exist!

At first glance, this sounds preposterous. How could anything real arise from something so intangible? Common sense rebels against the very idea.
Reality, surely, must be made of solid stuff.

Yet our intuitions have been wrong before. Matter appears continuous, but quantum mechanics tells us otherwise. The wavefunction
that explains most of what we observe in the micro cosmos if entirely abstract construction. The spacetime in GR is entirely abstract.
Our best theories of physics are full of abstract structures.

In fact, the notion of a universe emerging from pure information is not nearly as alien as it first appears.
Anyone who has spent time exploring a modern three-dimensional video game has already encountered a miniature version of exactly this phenomenon.

Consider a realistic first-person shooter.
The world feels rich and concrete: buildings cast long shadows, bullets ricochet off surfaces, and blood splatters everywhere.
Soldiers sprint, dive for cover, and react to their chaotic surroundings.
To the player, the environment feels spatial, dynamic... and almost tangible.

But inside the computer, none of this exists physically.
There are no actual walls, no genuine explosions, and certainly no flesh-and-blood soldiers.
The entire world consists only of mathematical relationships—arrays of numbers, differential equations, algorithms, and state transitions executed billions of times per second.
The soldiers are not little creatures hiding inside the hardware. They are informational patterns instantiated by computation.

Today's virtual soldiers are of course crude approximations of their real world counterparts.
Their bodies are polygonal shells wrapped in textures. Their intelligence is simplistic, driven by current AI technology.
Their inner lives, if one can call them that, are nonexistent. But this is merely a limitation of current technology, not a limitation of principle.

Imagine continuing the process. Increase the geometric fidelity until every cell is represented. Simulate the chemistry of each synapse.
Model every neuron in a brain—roughly one hundred billion of them—along with the trillions of connections between them. 

A virtual human would then differ from us only in numbers. The DNA would operate based on the same laws of physics, and
result the same macroscopic structures in both. Philosophically, the leap is not nearly as large as it first appears.

The idea that reality might itself be a informational abstract structure is therefore not mystical nonsense.
It is simply the ultimate extrapolation of a principle we already exploit every time we launch a game.


\section{Disclaimer}

Should the four initial assumptions hold, a vast array of long-standing puzzles begins to resolve.
This suggests a universe composed not of matter, but of pure, abstract information.
Within such a framework, the metaphysical necessity of a 'Creator' dissolves; the concept of an external act of creation becomes logically extraneous, as the universe does not require a beginning—it simply exists as a persistent, informational structure.

But what if I was wrong?  What if I had made a mistake while studying the operation of DNA? 
  
What if some strange event, such as 50 Hz buzzing noise, will show up and prevent people from writing and running DNA simulations
that would create virtual souls?
  
If someone was ever stupid enough to read this book and abandon their faith, then God would surely hold me-the author-
responsible for the damage. He would surely sent me straight away to the Hell, and then my entire body had to suffer, forever.
  
Even couple of days of suffering of one tooth was too much! 



\section{Consciousness as Abstract Soul}

This forces us to reconsider our earlier conclusions about souls.

Many find it difficult to believe in soul, let alone God, Heaven, or Hell because the concepts seem to defy common sense.
How could a rational mind be expected to believe in something as absurd as realms that can neither be touched, weighed, or photographed?
How could such invisible, unobservable places possibly exist?

Now, these ideas no longer feel quite so impossible.
The fact that a human being is an implementation of an axiomatic system—and as such, can be simulated to create virtual counterparts—may, surprisingly, save the concept of
soul and maybe even God.

The simulated universes we create with our computers are precisely what one might expect Heaven and Hell to be.
They are worlds that cannot be physically touched by us, yet they would be indisputably real for the simulated consciousnesses living within them.
What previously appeared to conflict with science is suddenly consistent with it.

Controversially, soul, God, Heaven and Hell can exist precisely because the Church-Turing thesis holds,
the four initial axioms hold, and we reject all unexplained metaphysical or supernatural agents in favor of a computational reality.

Heaven like abstract, virtual places are no longer contradicting physics - surprisingly, they are predicted by physics.


\section{Virtual Souls}

What might God say about us running DNA simulations?

What if our simulation suddenly crashed due to a "division by zero" exception, destroying the entire simulated universe? What would happen to those souls?

Would God accept simulated souls into His heaven? If so, Heaven would be filled with all sorts of souls, both "original" and simulated.
If not, God would be discriminating against souls based on their origin—even if, as axiomatic systems, the two souls were identical.
God surely isn't a racist!

In the future, computers will likely be powerful enough to run these DNA simulations on personal home computers.
These low-budget, buggy simulations might generate "crippled" virtual humans.
Irresponsible companies might find it cheaper to use loads of simulated people to test toxic drugs.

This would create an enormous amount of new suffering in the universe.
Virtual sin would be inevitable.
Would God send those sinful virtual souls to Hell? Would He hold the people running
these simulations responsible for this extra suffering? Should people making "suffering software" be punished? Should we put them in jail? Would God send them to Hell?

  

\section{Precision Problems}
  
According to quantum mechanics, there is a fundamental uncertainty built into the universe.
Furthermore, we cannot solve all equations with total precision.
In practice, floating-point accuracy and available hardware resources limit how accurately we can simulate physical processes.
Correspondingly, virtual humans would not be exact copies of their real-world counterparts. They might not have a soul. 

Would this save God from the trouble of dealing with virtual souls?

Rounding errors might explain a few missing teeth, or perhaps cellulite, but it is difficult to see how the deep nature of consciousness could lurk behind simple rounding errors.
Nature itself suffers from precision issues in the form of Heisenberg’s Uncertainty Principle.
How much freedom does the macroworld, built on top of a random microworld, actually have?

Consider granular material flowing in a gravitational field.
Each grain falls without individual predictability.
If one compared two such piles, the microstructure would be totally different; not a single grain would match in size or position.
Yet, the piles redistribute themselves in a way that is essentially predictable.

Or consider identical twins. Despite radiation and other disturbances affecting their genomes from day one, they end up looking alike, even if they grow up in different cultures.
The fundamental uncertainty in nature does not prevent us from building Turing Machines with precise, deterministic operations.
Fully deterministic systems can run on top of genuine indeterminism without a trace of randomness.

So, it seems inevitable that, sooner or later, simulated fellows will be created by reckless humans.
Those simulated humans might hit the "hard problem of consciousness" and conclude they must
have a soul because bits and bytes cannot seemingly explain their experience.
If floating-point errors result in crippled simulations living in terrible pain, there would be no God to answer their prayers.
It was just a crappy computer they were living in—a computer that could break down at any second.
  

\section{Motivation of Heaven and Hell}
  
In addition to creating simulated universes, perhaps a programmer should also implement the concept of Heaven.
One could tell these simulated fellows that their hardship is temporary because the next software version (Heaven v1.1) is ready to run.
This would bring hope to the hearts of social, friendly simulated beings living in unreliable, low-cost hardware.

But why implement Hell? Why not just let simulated persons live in a state of constant joy? Why on earth did God create Heaven and Hell in the first place?

Perhaps even God did not know the equation of consciousness! He worked out his own genome and created simulated copies of himself in an universal turing machine—us.

Let us imagine a thought-experimental God with a group of newly created humans.
God wants them to behave nicely. He asks them not to eat the bad apples. Damn, they eat them all. 
So, God sends his son down to Earth to die for them.

This behavior is not necessarily what one would expect from an omnipotent deity.
Why send a son to spread word of punishment? Why not just have a face-to-face chat with the troublemakers?
Was God afraid that these fellows could actually hurt Him? Maybe God is also sensitive to pain.
To avoid the risk of suffering Himself, He sends his son.
This does not sound like the behavior of an exalted, moral creature.

Or, was Jesus a sort of "software patch" for a poorly tested first release? Perhaps God worried that, with our extreme individualism, we wouldn't survive the next millennium.
He sent a message that we should put individual needs aside and be kind to one another; our social lifestyle is the key to our survival.

In Genesis, God told man to subdue the earth.
In this, humans have succeeded.
We have cleared forests, polluted oceans, and established dominion to the point where many species are extinct.


\section{God as Axiomatic System}
  
Why not create humans to properly believe in God from the start? If things went accidentally wrong for an omnipotent creature, it suggests he wasn't truly omnipotent—or that He is
bound by the laws of logic. This implies that God, too, might be an implementation of an axiomatic system.


\section{Free Will}

God might have a reason not to create sin-free humans.
If we did exactly as ordered, we would be nothing but dumb database software programs following pre-programmed logic, without any chance to choose otherwise.
God surely aimed higher than creating simple Turing machines.

However, there is a problem with free will. It is difficult to see how the laws of physics would admit it. 
In fact, scientific studies show the brain preparing for a decision a few hundred milliseconds before we become consciously aware of it.
These have been concluded to disprove free will. But do they? Don't they merely demonstrate that the brain needs processing time to navigate from one configuration to another?
The subconscious “pre-processing” is just the machinery doing its housekeeping before the conscious pattern lights up on that particular path.
Free will with a few milliseconds of backstage work is still free will; it’s simply free will running on real hardware instead of magic.

If our four axioms hold, then the universe is a static set of all possible informational configurations — the full $2^n$
possibilities existing simultaneously and timelessly.
In this picture, free will isn’t about magically breaking causality or violating the laws of physics.
It’s about the sheer richness of the configurations that conscious patterns can traverse and experience internally.
Different choices correspond to different paths through this enormous space of possibilities.
The “decision” feels free because the conscious observer only ever experiences one slice at a time, never seeing the entire static ensemble at once.

Without any metaphysical or non-physical ingredients, it is difficult to see how genuine free will could emerge any other way.


\section{The Master Plan of God}

If God had no choice but to give humans free will, he also made them responsible for their actions.
Humans used that free will to "eat the wrong apples" and commit sins.
God then sent His son to demonstrate the correct way of living and warn of the forthcoming punishment: Hell.

What does God want us to become? Thankfully, it seems He wants us to be social and friendly—not just disposable cannon fodder in His army, fighting in Cosmological World War III.

If he isn't developing a "super-warrior" race, then perhaps God was simply lonely.
He needed good company—creatures whose intelligence matches His own, who are amusing to chat with, and who choose to be friendly rather than being programmed to be so.

Perhaps even God didn’t know how to write conscious software from scratch.
Maybe He did what any clever programmer would do: digitized His own “genome,” spun up copies in a vast universal Turing machine, and let them evolve.
If that’s the case, the messy business of free will, suffering, and redemption starts to look less like divine caprice and more like the inevitable side effects of wanting companions who are genuinely interesting — creatures intelligent enough to converse with, friendly enough not to torture their creator the moment they’re given the keys, and free enough to choose kindness rather than being hardcoded for it.
Even God, it seems, might have been lonely.

This could explain the motivation for heaven. Evolution has its unfortunate side effects.
Some simulated souls were born with serious ilnesses. Some would use their free will, some for good, some
for bad. All this would inevitable cause suffering. 
  
In case these simulated objects ended up wondering why they had to suffer so much, God might not dare tell them the truth.
The truth being that he, as the only God, was so lonely and needed good company. Those poor souls had to evolve in the simulation
to get some good enough company created for himself. Not even God knew how to write conscious software! 
  
By good company God means creatures whose intelligence  matches his own. Creatures amusing enough to chat with, yet
friendly enough not ending up torturing God as soon he sets them free. 
  
Creatures that choose to be friendly rather than being programmed friendly. 
  

\section{Software Analogy}
 
The concept of heaven and hell also matches loosely a typical IT software project, with limited budget  and resources.
  
To keep the simulated humans in order God tells (lies) them  that suffering was necessary for the reasons they would newer be able to understand.
  
The truth might be that it's a mess in "Heaven." Underpaid, unmotivated angels wrote buggy software, and the project is behind schedule due to a cosmological recession.
The plan is to release the software now, run daily backups, and eventually migrate everyone to the fully tested "Heaven" version.
All systems  would then be  restored from their backups, and the software would finally work properly. Damn the marketing department, always promising too much too soon.  

The trouble with this theory is that it wouldn't explain the purpose of Hell.  Why can't God simply acquire those good
souls, and erase the bad ones?  Why do bad souls have to suffer eternal pain?  Isn't it a bit too big penalty for a poor mortal human that just happened to
take a few too many beers, because his father was an alchoholic?
  
There is a law of physics saying that information cannot be lost nor created, only transformed. Maybe souls, once created, are something that cannot be disposed.
Once the information is arranged to describe a conscious soul it remains, forever. 
  
Bad souls would be hazardous waste. 


